% SITWE 2026 — Symposium on Innovations and Technologies in Web Engineering
% Tan Poi Cheong, Multimedia University
% 6-page IEEE conference paper, IEEEtran two-column, 10 pt Times.
%
% Compile with:  pdflatex sitwe2026_paper.tex  (twice for cross-refs)
% or open in Overleaf with IEEEtran class enabled.

\documentclass[conference,10pt,a4paper]{IEEEtran}

\IEEEoverridecommandlockouts

\usepackage{cite}
\usepackage{amsmath,amssymb}
\usepackage{graphicx}
\usepackage{booktabs}
\usepackage{array}
\usepackage{url}
\usepackage{hyperref}
\usepackage{xcolor}
\usepackage{textcomp}
\usepackage{listings}
\usepackage{tikz}
\usetikzlibrary{positioning, arrows.meta, shapes.geometric, fit, backgrounds}

\lstset{
  basicstyle=\ttfamily\footnotesize,
  breaklines=true,
  columns=fullflexible,
  keepspaces=true,
  showstringspaces=false,
}

\begin{document}

\title{A Schema-Aware Model Context Protocol Server for Natural Language to
SQL Translation under Schema Evolution\\
\thanks{This research is conducted at the Faculty of Computing and
Informatics, Multimedia University, Cyberjaya, Malaysia, as part of an
M.Sc.\ (IT) by Research programme.}}

\author{\IEEEauthorblockN{Poi-Cheong Tan}
\IEEEauthorblockA{Faculty of Computing and Informatics\\
Multimedia University\\
Persiaran Multimedia, 63100 Cyberjaya, Malaysia\\
tan.poi.cheong@student.mmu.edu.my}
\and
\IEEEauthorblockN{Su-Cheng Haw}
\IEEEauthorblockA{Faculty of Computing and Informatics\\
Multimedia University\\
Persiaran Multimedia, 63100 Cyberjaya, Malaysia\\
sucheng@mmu.edu.my}
\and
\IEEEauthorblockN{Hui-Ngo Goh}
\IEEEauthorblockA{Faculty of Computing and Informatics\\
Multimedia University\\
Persiaran Multimedia, 63100 Cyberjaya, Malaysia\\
hngoh@mmu.edu.my}
\and
\IEEEauthorblockN{Jayapradha J}
\IEEEauthorblockA{Department of Computing Technologies,\\
School of Computing\\
SRM Institute of Science and Technology\\
Kattankulathur, Tamil Nadu 603203, India\\
jayapraj@srmist.edu.in}
}

\maketitle

\begin{abstract}
Natural Language to SQL (NL-to-SQL) systems have advanced rapidly,
yet every mainstream system is trained and evaluated against a fixed,
pre-defined schema. In production, schemas evolve: EvoSchema shows state-of-the-art
systems including GPT-4 degrading sharply under table- and
column-level perturbations. The
Model Context Protocol (MCP) lets an LLM client discover schemas
and invoke query tools at runtime, yet no peer-reviewed study has
evaluated it on a standard NL-to-SQL benchmark. This paper closes that gap with a schema-aware MCP server
exposing three runtime primitives ---
\texttt{schema/fingerprint} (BLAKE3 hashing with diff classification),
\texttt{query/relink} (AST rewriting with LLM fallback), and
\texttt{query/validate} (execution feedback with back-translation) ---
and an open-source pilot harness.
On a 76-query evaluation over a perturbed Spider \texttt{concert\_singer}
database spanning all five EvoSchema operator classes with Claude
Haiku~4.5, the MCP-mediated path reaches a Recovery Rate of $0.80$.
$\text{EX}$ rises from a $0.11$ stale-schema baseline to $0.59$,
against a $0.71$ refreshed-schema ceiling.
Across four LLM clients and two Spider sub-domains, Recovery Rates
range from $0.80$ to $0.97$, consistent with the server's
model-agnostic design.
A like-for-like vendor-native function-calling baseline is
statistically indistinguishable from the MCP path, and on a 199-item
EvoSchema/BIRD-dev subset the server restores execution accuracy from
a near-zero stale floor to the refreshed-schema ceiling
(RR $= 1.08$--$1.15$, $p < 10^{-6}$).
Residual non-recovery concentrates on column-merge and a
value-encoding mismatch (Q05). The server, harness, and evaluation
dataset are released under the MIT licence.
\end{abstract}

\begin{IEEEkeywords}
Natural Language to SQL, Schema Evolution, Model Context Protocol,
Large Language Models, Database Systems.
\end{IEEEkeywords}

% --------------------------------------------------------------- %
\section{Introduction}
% --------------------------------------------------------------- %

Translating natural-language questions into executable SQL has progressed
through four distinct eras: rule-based parsers (pre-2017), sequence-to-sequence
neural parsers such as Seq2SQL~\cite{zhong2017seq2sql} and IRNet, graph-augmented
Transformers including LGESQL and RAT-SQL with constrained decoding via
PICARD, and most recently decoupled pipelines such as
RESDSQL~\cite{li2023resdsql} and LLM-prompting frameworks
DIN-SQL~\cite{pourreza2023dinsql}, DAIL-SQL, and multi-agent systems
including MAC-SQL~\cite{wang2025macsql}, CHASE-SQL~\cite{pourreza2025chasesql}, and CHESS~\cite{talaei2024chess}. Across this
ten-year trajectory, one experimental condition has remained constant: a
fixed database schema is fully known at inference time. Benchmark corpora
including Spider~\cite{yu2018spider}, BIRD~\cite{li2023bird}, and Spider~2.0
all enforce this assumption.

In production environments this assumption is systematically violated.
Brahmia et al.~\cite{brahmia2024schemaevol} survey four decades of
schema-evolution research in database systems and identify five operator
classes --- table rename, table split, table merge, column rename, column
merge --- that occur routinely after deployment. The recent
EvoSchema~\cite{zhang2025evoschema} benchmark provides the first large-scale
empirical evidence of degradation: applying ten perturbation types to Spider
and BIRD, every tested model — including GPT-4 — suffers significant
accuracy drops, with table-level changes substantially more damaging than
column-level changes. EvoSchema, however, measures the problem without
proposing a runtime remedy.

A parallel development supplies the mechanism that statically prompted
NL-to-SQL systems lack. The Model Context Protocol
(MCP)~\cite{anthropic2024mcp}, introduced by Anthropic in November 2024 and
surveyed comprehensively by Hou et al.~\cite{hou2025mcp}, defines a
session-oriented JSON-RPC protocol through which an LLM can dynamically
discover \emph{tools}, \emph{resources}, and \emph{prompts} exposed by
external systems. Database connectors are among the dominant categories of
MCP server deployed today; early proof-of-concept systems
Toolbelt-MCP~\cite{agater2026toolbelt} and EHR-MCP~\cite{masayoshi2025ehrmcp}
demonstrate that an LLM client can discover schema metadata and issue SQL
through MCP without any fine-tuning. Neither system has been evaluated on
Spider, BIRD, or EvoSchema, and neither addresses runtime schema evolution.

\paragraph*{Research gap}
A systematic literature scan (first conducted May~2026, refreshed
July~2026) confirms that no published study isolates MCP-mediated
dynamic schema retrieval as an experimental variable on a standard
NL-to-SQL benchmark, and no system detects, classifies, or recovers
from schema change through the protocol. The nearest work,
AgentNLQ~\cite{bogdanov2026agentnlq}, employs MCP only as tool
transport over a schema that is static within and across sessions and
reports no schema-change condition. This absence defines the gap
addressed here.

\paragraph*{Contributions}
This paper makes the following contributions:
\begin{enumerate}
  \item A schema-aware MCP server exposing three runtime primitives ---
    \texttt{schema/fingerprint}, \texttt{query/relink}, and
    \texttt{query/validate} --- that together enable an off-the-shelf LLM
    client to detect, represent, and recover from schema evolution at
    inference time without fine-tuning.
  \item A 76-query pooled evaluation on perturbed Spider
    \texttt{concert\_singer} data spanning all five EvoSchema operator
    classes that empirically isolates schema retrieval as the
    independent variable and reports Execution Accuracy and Recovery
    Rate, and a matching 76-query evaluation on a second sub-domain
    (\texttt{hr\_1}, Section~\ref{sec:results}-C), for 152 queries in
    total.
  \item An open-source reference implementation (the server, evaluation
    harness, and pilot fixture) released to lower the cost of subsequent
    research on MCP-mediated database access. The code, the pilot
    dataset, the perturbed SQLite fixtures, and the figure-rendering
    scripts are released at \url{https://github.com/pctan88/sitwe2026-mcp-nlsql}
    under the MIT
    licence. Recorded versions: \texttt{sqlglot 23.x}, \texttt{blake3-py
    0.4.x}, \texttt{python-docx 1.1.x}, Python~3.11.
\end{enumerate}

% --------------------------------------------------------------- %
\section{Related Work}
% --------------------------------------------------------------- %

\subsection{NL-to-SQL Systems and Benchmarks}
Recent surveys contextualise the state of LLM-era NL-to-SQL.
Shi et al.~\cite{shi2025survey} survey LLM-based text-to-SQL methods
in \emph{ACM Computing Surveys}, covering prompt-engineering and
fine-tuning strategies and establishing no-fine-tuning prompting as a
recognised research path. Hong et al.~\cite{hong2025tkde}
provide a complementary survey tracing the shift from pre-trained to
large language models, and identify schema linking and
execution-grounded feedback as the two sub-tasks on which most
accuracy gains concentrate.

Among recent high-performing systems, CHASE-SQL~\cite{pourreza2025chasesql}
achieves state-of-the-art execution accuracy on BIRD through multi-path
reasoning and preference-optimised candidate selection;
MAC-SQL~\cite{wang2025macsql} employs a multi-agent collaborative
framework with specialised decomposer, refiner, and selector agents.
DeepEye-SQL~\cite{li2026deepeyesql} adopts a software-engineering-inspired
modular pipeline reported at ACM SIGMOD 2026.
Lei et al.~\cite{lei2026crl} decompose the task into schema linking and
SQL generation, training each stage with a dedicated cooperative
reinforcement-learning reward to improve robustness on complex
multi-table joins.
None of these systems addresses runtime schema evolution.

Two peer-reviewed studies motivate the enterprise importance of the
problem. Ojuri et al.~\cite{ojuri2025agents} combine GPT-based LLMs
with intelligent agents for enterprise NL-to-SQL,
demonstrating that a ReAct agent architecture outperforms few-shot
in-context learning for business-analytics tasks, yet still assumes a
static schema at inference time. Nascimento et
al.~\cite{nascimento2026kgsql} address real-world schema challenges,
using a knowledge graph to expose schema semantics to an LLM together with
dynamic few-shot example selection; their analysis confirms that Spider
and BIRD do not capture the schema-complexity issues present in
production databases, directly motivating the runtime schema-fingerprint
primitive proposed in this paper.

\subsection{NL-to-SQL under Schema Variation}
Liao et al.~\cite{liao2026infosql} formalise an adversarial-defence task
spanning perturbations to the question, the schema, and the SQL logic,
and report state-of-the-art robustness across five Spider-derived
benchmarks; their setting, however, is addressed through model training
rather than a runtime protocol.
CHESS~\cite{talaei2024chess} achieves 71.10\,\% on
BIRD with 83\,\% fewer LLM calls; its self-consistency check over
multiple candidates catches some silent failures, but it does not
employ SQL-to-text back-translation against the original NL question
as a dedicated silent-failure signal.
GBV-SQL~\cite{chen2025gbvsql} introduces SQL-to-text back-translation as
a self-check mechanism but operates only on static schemas and is not
exposed as a standardised protocol.
Li et al.~\cite{li2025cmsql} propose a cross-model consistency framework
that validates SQL candidates by combining fine-tuned models with LLM
reasoning; their consistency check, however, operates within a
single vendor's pipeline rather than over an open protocol.
TRUST-SQL~\cite{jian2026trustsql}
formalises the ``Unknown Schema'' setting closest to our problem framing
but requires multi-turn reinforcement-learning training out of scope of
the present work.

\subsection{Schema Evolution in Database Systems}
Brahmia et al.~\cite{brahmia2024schemaevol} compare schema-evolution
proposals along six criteria (database model, implementation, change
semantics, propagation, integrity, software evolution) and provide the
most recent comprehensive review of the area. Their taxonomy of five operator types
(table rename, table split, table merge, column rename, column merge)
defines the diff classification implemented by our
\texttt{schema/fingerprint} Resource.
At the system level, recent work addresses complementary
schema-related obstacles in production settings: Lei et
al.~\cite{lei2026srspsql} combine question semantic reconstruction with
dynamic schema pruning to reduce ambiguity and schema redundancy,
demonstrating that schema-aware pre-processing measurably improves SQL
generation on Spider and BIRD.
Such methods assume the schema is fixed at inference time, whereas the
\texttt{schema/fingerprint} Resource proposed here treats schema change
itself as the runtime signal.

\subsection{Model Context Protocol and Tool-Augmented LLMs}
MCP generalises pre-existing tool-use patterns
(Toolformer, ReAct, Gorilla) into an open
standard. Hou et al.~\cite{hou2025mcp} characterise MCP as an
interoperability layer for AI agents, defining a threat taxonomy
spanning four attacker categories across the MCP lifecycle,
with tool poisoning as a prominent risk. Toolbelt-MCP~\cite{agater2026toolbelt} demonstrates
schema discovery and SQL execution against a $\sim$1\,000-table industrial
database; EHR-MCP~\cite{masayoshi2025ehrmcp} demonstrates the same pattern
over a Japanese teaching-hospital electronic health record. Neither system
has been evaluated on Spider or BIRD, and neither addresses schema evolution.
More recently, AgentNLQ~\cite{bogdanov2026agentnlq} reports 78.1\,\%
semantic accuracy on BIRD through a multi-agent orchestrator that uses MCP
as tool transport, but its schema enrichment is prepared offline and no
schema-change condition is evaluated. DynaQuery~\cite{hassini2025dynaquery}
elevates live schema introspection to a query-planning phase, but neither
classifies evolution events against a taxonomy nor measures recovery, and
does not use MCP. Schema-change awareness through the protocol therefore
remains unclaimed.

% --------------------------------------------------------------- %
\section{Methodology}
% --------------------------------------------------------------- %

The system is organised in three layers separated by clear responsibility boundaries. Fig.~\ref{fig:arch} shows the MCP-mediated NL-to-SQL query cycle: the user submits a natural-language query (step~1); the LLM client interprets it and invokes the MCP server (step~2); steps~3--5 (shaded region) are handled server-side via \texttt{schema/fingerprint}, \texttt{query/relink}, and \texttt{query/validate}; step~6 returns the validated response to the user.

\begin{figure}[t]
\centering
\includegraphics[width=\columnwidth]{paper/figures/fig_arch_cycle.png}
\caption{The MCP-mediated NL-to-SQL query cycle.}
\label{fig:arch}
\end{figure}

Layer~1 receives a natural-language question from the user. Layer~2 hosts a
pre-trained, MCP-compatible LLM client (Claude Haiku~4.5 in the pilot
described below); the client is used as released by the vendor with no
parameter updates of any kind. Layer~3 is the MCP server proposed by this
work, exposing the three primitives detailed below. Layer~4 holds the
SQLite databases distributed by Spider, BIRD, and EvoSchema, against which
candidate SQL is executed read-only.

\subsection{Schema Fingerprint Resource}
Following the Apache Avro Schema-Fingerprints precedent, the
\texttt{schema/} \texttt{fingerprint} Resource computes a canonical hash
\begin{equation}
\mathcal{F}(S) \;=\; H\!\left( \bigsqcup_{T \in S}\!\!\left[ \text{name}(T)
\,\|\, \bigsqcup_{C \in T}\!\!\left( \text{name}(C) \,\|\, \text{type}(C)
\right) \right] \right)
\label{eq:fp}
\end{equation}
where $H$ is the BLAKE3 hash function ($256$-bit digest); the operator
$\bigsqcup$ denotes \texttt{concat\_sorted} --- byte concatenation of
its operands after they are sorted lexicographically (applied at both
the table-set and column-set levels, making the result invariant under
table reordering and column reordering); and $\|$ is plain byte
concatenation within a single tuple. We use $\bigsqcup$ rather than
$\bigoplus$ to avoid the XOR connotation that the latter carries in
the cryptographic literature. Two consecutive
fingerprints are equal --- with cryptographic probability --- iff the
schema content (table names, column names, column types) is unchanged; BLAKE3's collision resistance bounds the
probability of a coincidental match below $2^{-128}$. When they
differ, a per-object diff is computed and classified into one of the five
operator types of Brahmia et al.~\cite{brahmia2024schemaevol}: TABLE\_RENAME,
TABLE\_SPLIT, TABLE\_MERGE, COLUMN\_RENAME, COLUMN\_MERGE. The Resource
returns the schema in compact column-list format with token-budget
instrumentation, following the empirical observation in
EHR-MCP~\cite{masayoshi2025ehrmcp} that long resource responses risk LLM
context-window overflow.

\subsection{Query Relink Tool}
The \texttt{query/relink} Tool accepts a stale SQL string and a classified
diff, and returns rewritten SQL targeting the evolved schema. The four
straightforward operator types (table rename, table merge, column rename,
column merge) are rewritten using deterministic AST transformations via the
\texttt{sqlglot} library. The two complex operator types
(TABLE\_SPLIT and COLUMN\_MERGE-with-expression) trigger an LLM-
assisted fallback that takes the stale SQL,
the diff, and the new schema as a single prompt and returns the rewritten
query. Because the AST path is deterministic, repeated runs over the same
inputs produce identical outputs.

\subsection{Query Validate Tool}
The \texttt{query/validate} Tool executes candidate SQL against the live
database and returns one of three verdicts: \texttt{valid},
\texttt{execution\_error}, or \texttt{silent\_failure\_suspected}. The silent-
failure heuristic suite combines (i)~result-set arity checks, (ii)~empty
result on an affirmative question, and (iii)~SQL-to-text back-translation
into a paraphrase that is compared to the original question through a
token-overlap metric. The v2 detector realises this with a local
MiniLM sentence-embedding cosine check
(Section~\ref{sec:results}). The validation step is performed
entirely inside the MCP server, keeping the LLM-client cost bounded to one
generation pass plus one optional re-prompt after a validation failure.

\subsection{Protocol Surface}
All three primitives are exposed over MCP's JSON-RPC~2.0 transport, making
the server usable from any MCP-compatible client without
vendor-specific function-calling adapters. This is the qualitative property
that distinguishes the proposed work from vendor-coupled agent
architectures such as AskDB~\cite{phan2025askdb} and
CHESS~\cite{talaei2024chess}.

% --------------------------------------------------------------- %
\section{Pilot Study Setup}
% --------------------------------------------------------------- %

\subsection{Database and Perturbations}
A self-contained \texttt{concert\_singer} SQLite database — mirroring the
Spider~1.0 schema for that domain — was created in two versions: a
pre-perturbation copy (tables \texttt{stadium}, \texttt{singer},
\texttt{concert}, \texttt{singer\_in\_concert}) and a post-perturbation
copy in which the table \texttt{singer} is renamed to \texttt{artist}, with
\texttt{Song\_release\_year} renamed to \texttt{debut\_year} and
\texttt{Song\_Name} to \texttt{song\_title}. These perturbations correspond
to one TABLE\_RENAME and two COLUMN\_RENAME operators from the
EvoSchema taxonomy.

\subsection{Pilot Query Set}
Twenty natural-language questions were authored against the
\texttt{concert\_singer} domain. Each question carries (i)~a gold SQL
answer against the pre-schema and (ii)~a gold SQL answer against the
post-schema. All forty gold queries execute successfully and return
equivalent result sets across the two databases. The 20 questions
decompose as: eleven exercise only the TABLE\_RENAME operator, seven
exercise the combined TABLE\_RENAME plus COLUMN\_RENAME operators, and
two are control queries that touch only the unmodified \texttt{stadium} or
\texttt{concert} tables.

\subsection{Configurations}
Four configurations are evaluated, all sharing the same stale-schema
prompt so that runtime schema retrieval is the only variable that changes
across them:
\begin{description}
  \item[Baseline (stale schema)] The LLM client receives the \emph{stale}
    (pre-perturbation) schema as part of its prompt; the generated SQL is
    executed against the \emph{post-perturbation} database. This mirrors
    the prevailing paradigm of static-schema prompting.
  \item[Refreshed-schema] The LLM client is given the post-perturbation
    schema directly in its prompt and regenerates the SQL from scratch,
    with no MCP server in the loop. This path measures the accuracy
    ceiling attainable when the model simply sees the current schema, and
    it serves as the reference against which recovery is scored.
  \item[Error-feedback] The LLM client receives the stale-schema prompt;
    when the generated SQL raises an execution error on the
    post-perturbation database, the error message is returned to the model
    for a single retry, still under the stale schema. This isolates the
    contribution of a naive one-shot error-correction loop that lacks any
    schema diff.
  \item[MCP-mediated] The LLM client receives the same stale-schema prompt.
    The generated SQL is then passed through the MCP server: the
    \texttt{schema/fingerprint} Resource computes the diff against the live
    schema; \texttt{query/relink} rewrites the SQL using AST transformations
    driven by the diff; \texttt{query/validate} executes the rewritten SQL
    and emits a verdict, triggering a single LLM re-prompt on a
    \texttt{silent\_failure\_suspected} verdict.
\end{description}
The LLM client is held constant across all four configurations to isolate
runtime schema retrieval as the independent variable. All results
reported here use Claude Haiku~4.5 (model identifier
\texttt{claude-haiku-4-5-20251001}) via the Anthropic Messages API with
temperature fixed at zero. Note that \texttt{temperature = 0} on the
Anthropic API reduces but does not eliminate non-determinism across
calls; the canonical pilot records a single completion per query. The
harness's \texttt{--seeds N} flag re-runs each query $N$ times at
temperature $0$; a five-seed sweep confirmed
\texttt{seed\_stability\_baseline} $= 1.0$
(Section~\ref{sec:results}-C). A deterministic NL-to-SQL stub
supports unit testing without API cost.

\subsection{Metrics}
Execution Accuracy (EX) is computed as the fraction of questions for which
the predicted SQL produces the same result set (treated as a row-order-
independent multiset) as the gold SQL. Recovery Rate (RR) is defined
against the refreshed-schema ceiling as
\begin{equation}
RR \;=\; \frac{\text{EX}_{\text{mcp}} - \text{EX}_{\text{stale}}}
              {\text{EX}_{\text{refreshed}} - \text{EX}_{\text{stale}}}
\label{eq:rr}
\end{equation}
where $\text{EX}_{\text{stale}}$ is the stale-schema baseline accuracy and
$\text{EX}_{\text{refreshed}}$ is the accuracy of the same LLM when it is
given the post-perturbation schema directly (the Refreshed-schema
configuration). $RR = 1.0$ denotes full recovery to the refreshed-schema
ceiling; $RR = 0$ denotes no recovery; $RR < 0$ denotes degradation; values
marginally above $1.0$ can arise when relinking a previously correct query
outperforms regeneration from scratch. The denominator uses
$\text{EX}_{\text{refreshed}}$ rather than the unperturbed-schema reference
$\text{EX}_{\text{pre}}$ so that recovery is scored against what the model
could achieve with full knowledge of the evolved schema. Silent-Failure
Detection Rate is reported as the fraction of baseline failures correctly
flagged by \texttt{query/validate} on the MCP path.

% --------------------------------------------------------------- %
\section{Preliminary Results}
\label{sec:results}
% --------------------------------------------------------------- %


Table~\ref{tab:results} summarises the pooled results across all five
EvoSchema operators ($n = 76$ queries). The substantive empirical
content is three falsifiable propositions: (i)~the diff classifier
identifies all ground-truth perturbations with zero false positives on
the unperturbed control queries; (ii)~the \texttt{query/relink} Tool
produces parseable, executable SQL for every recoverable query, and
the \texttt{query/validate} Tool does not suppress any correct rewrite
via a false-positive silent-failure verdict; and (iii)~the end-to-end
MCP transport composes without protocol failure across all sessions.
The headline quantitative result is: $\text{EX}$ raised from $0.11$
under the stale-schema baseline to $0.59$ under the MCP-mediated
path, against a refreshed-schema ceiling of $0.71$, giving Recovery
Rate $0.80$ by Equation~\ref{eq:rr}. Residual non-recovery ($0.20$)
concentrates on the column-merge operator --- where inverse-expression
reconstruction lies outside the AST rewriter's scope --- and on Q05
(analysed below), which shows that identifier-only AST rewriting
cannot normalise LLM-introduced value-level predicate mismatches.

The stale-schema baseline returned $\text{EX} = 0.11$ across all $76$
queries ($8/76$ correct). Execution errors dominate: $63$ of the $76$
baseline queries fail outright because the stale SQL references
renamed or restructured objects that no longer exist in the evolved
schema (see Table~\ref{tab:erroranalysis}). The eight correct baseline
queries arise where a stale identifier happens to resolve against a
surviving schema object. The five wrong-result cases occur where a
stale identifier resolves against a surviving object but yields an
incorrect result set. The Wilson 95\,\% CI of $[0.05, 0.19]$ on the
baseline is reported in Table~\ref{tab:results}.

\begin{table}[h!]
\centering
\caption{Execution Accuracy and Recovery Rate across the four
configurations, pooled over all five EvoSchema operators on the perturbed
\texttt{concert\_singer} database with Claude Haiku~4.5 ($n=76$). EX shown
with Wilson 95\,\% CI; RR by Equation~\ref{eq:rr}. ``$\Delta$ vs B'' is the
EX gain over the stale-schema baseline.}
\label{tab:results}
\begin{tabular}{lccc}
\toprule
\textbf{Configuration} & \textbf{EX [95\,\% CI]} & $\Delta$ \textbf{vs B} & \textbf{RR} \\
\midrule
Baseline (stale schema)     & 0.11 [0.05, 0.19] & --              & --            \\
Refreshed-schema (ceiling)  & 0.71 [0.60, 0.80] & +0.61           & --            \\
Error-feedback (1 retry)    & 0.11 [0.05, 0.19] & +0.00           & --            \\
MCP-mediated (proposed)     & \textbf{0.59 [0.48, 0.70]} & \textbf{+0.49} & \textbf{0.80} \\
\bottomrule
\end{tabular}
\end{table}

Fig.~\ref{fig:ex_bar} plots Execution Accuracy for all
four configurations with Wilson 95\,\% confidence intervals;
the arrow marks the pooled Recovery Rate of $0.80$.

\begin{figure}[h!]
\centering
\includegraphics[width=\columnwidth]{paper/figures/fig_ex_bar.pdf}
\caption{Execution Accuracy across four configurations. Error bars: Wilson 95\,\% CI. Arrow: RR = 0.80.}
\label{fig:ex_bar}
\end{figure}

Table~\ref{tab:perpert} disaggregates results by operator across the
76-query evaluation. The TABLE\_RENAME, TABLE\_SPLIT, and COLUMN\_RENAME
operators show high recovery (MCP correct $\geq 90\%$ of pre-time
correct queries) because the deterministic AST relinker handles
identifier substitution without error. TABLE\_MERGE recovery is
slightly lower; notably the MCP path corrects $11$ queries against a
pre-perturbation ceiling of $10$, giving a per-operator RR of $1.125$.
This marginal overshoot occurs when relinking outperforms regeneration
from scratch (same effect as cross-model RR $> 1.0$ noted in
Section~\ref{sec:results}-C). The COLUMN\_MERGE operator is the
consistent exception: when two columns are merged through an expression,
the original components cannot be recovered from the post-schema; all
four models show near-zero incremental recovery on
\texttt{concert\_singer} COLUMN\_MERGE under the MCP path. Failures
attributable to the LLM generating incorrect SQL on the unperturbed
schema (pre-time competence misses) lie outside the scope of the
schema-evolution mechanism.
Stratifying all pilot queries by gold-SQL structure (easy: single
table; medium: one join or aggregation) shows the recovery is not an
artefact of trivial queries. Pooled across the four evaluated clients,
RR is $0.92$ on the easy bucket ($n{=}344$) and $0.93$ on the medium
bucket ($n{=}264$). The authored pilot sets contain no hard-bucket
queries ($\geq 2$ joins, subqueries, \texttt{HAVING}); harder
questions are exercised by the EvoSchema subset
(Section~\ref{sec:evoschema}).

\begin{table}[h!]
\centering
\caption{Per-perturbation breakdown across all five EvoSchema operators
($n = 76$, \texttt{concert\_singer}, Haiku~4.5). Pre = refreshed-schema
correct (recovery ceiling); Base = stale-schema baseline correct;
MCP = MCP-mediated correct; Rec.\ = queries recovered by MCP (MCP $-$ Base).}
\label{tab:perpert}
\begin{tabular}{lccccc}
\toprule
\textbf{Pert.} & \textbf{N} & \textbf{Pre} & \textbf{Base} & \textbf{MCP} & \textbf{Rec.} \\
\midrule
TABLE\_RENAME   & 16 & 12 & 1  & 9  & 8  \\
TABLE\_SPLIT    & 16 & 12 & 1  & 10 & 9  \\
TABLE\_MERGE    & 14 & 10 & 2  & 11 & 9  \\
COLUMN\_RENAME  & 16 & 13 & 1  & 12 & 11 \\
COLUMN\_MERGE   & 14 & 7  & 3  & 3  & 0  \\
\midrule
Total           & 76 & 54 & 8  & 45 & 37 \\
\bottomrule
\end{tabular}
\end{table}

\paragraph*{Q05: a value-encoding mismatch the rewriter cannot reach}
Q05 asks ``How many male singers are there?''; the gold answer is a
count over \texttt{Is\_male = 'T'}. Haiku hedged the boolean encoding
with a wide OR-list: \texttt{SELECT COUNT(*) FROM singer WHERE
Is\_male = 'true' OR Is\_male = 'True' OR Is\_male = 'T' OR Is\_male =
'yes' OR Is\_male = 'Yes' OR Is\_male = 'Y' OR Is\_male = '1'}. This
matched gold pre-time only because the data happens to contain
\texttt{'T'}. \texttt{query/relink} renamed \texttt{singer} $\to$
\texttt{artist} faithfully and kept every disjunct, so the post-schema
query runs without error and returns $0$ where gold returns the real
count: a silent semantic failure that survives a correct identifier
rewrite. The empty-result-on-affirmative-question heuristic flags it
(\texttt{COUNT(*) = 0} on a ``how many'' question), but the re-prompt
does not repair the encoding, so detection does not imply recovery for
value-level mismatches.

The \texttt{schema/fingerprint} Resource correctly identified one
schema-change event per operator-perturbed database copy, with zero
false positives on unperturbed copies, confirming BLAKE3 hashing
reliably separates evolved from unmodified schemas across all five
operator classes. For the four AST-tractable operators,
\texttt{query/relink} required no internal LLM fallback of its own.

The silent-failure detector is evaluated on a purpose-built labelled
set: $60$ cases ($20$ faithful, $40$ injected silent failures)
spanning all five operators on both pilot databases, with four
mutation families:
wrong-column substitution, dropped \texttt{WHERE} clause, aggregate
swap, and stale merge-equality predicates. Table~\ref{tab:sfd} reports
per-operator precision, recall, and F1 for two backends. v1 is the
token-overlap back-translator; v2 is a sentence-embedding cosine check
(MiniLM, local inference, threshold $0.45$ from a sweep). The v2
backend reaches precision $1.00$ with recall $0.375$ and F1 $0.545$:
no false positive on the $20$ faithful queries, so no correct rewrite
is suppressed, and it beats v1 (precision $0.92$, recall $0.30$) on
both axes. Detection is
strongest where the pilot needs it. The stale merge-equality mutation,
the Q05 value-encoding family, is caught $4/4$; aggregate swaps $3/3$.
Most misses are wrong-column and dropped-predicate mutations whose
back-translations read almost the same in English. Precision-first is deliberate: a false positive discards a correct
rewrite; a miss leaves the query no worse than the baseline.

\begin{table}[t]
\centering
\caption{Silent-failure detection on the 60-case labelled set
(both databases, threshold $0.45$). v2 = sentence-embedding cosine
backend; pooled v1 token-overlap baseline shown for comparison.
FP $= 0$ for every operator under v2.}
\label{tab:sfd}
\begin{tabular}{lccccc}
\toprule
\textbf{Operator} & \textbf{N} & \textbf{TP} & \textbf{Prec.} & \textbf{Rec.} & \textbf{F1} \\
\midrule
TABLE\_RENAME  & 12 & 3 & 1.00 & 0.375 & 0.545 \\
TABLE\_SPLIT   & 12 & 3 & 1.00 & 0.375 & 0.545 \\
TABLE\_MERGE   & 12 & 0 & ---  & 0.000 & 0.000 \\
COLUMN\_RENAME & 12 & 4 & 1.00 & 0.500 & 0.667 \\
COLUMN\_MERGE  & 12 & 5 & 1.00 & 0.625 & 0.769 \\
\midrule
Pooled (v2)    & 60 & 15 & 1.00 & 0.375 & 0.545 \\
Pooled (v1)    & 60 & 12 & 0.92 & 0.300 & 0.453 \\
\bottomrule
\end{tabular}
\end{table}

On a 2024 commodity laptop the mean MCP middleware overhead
(\texttt{schema/} \texttt{fingerprint} diff $+$ \texttt{query/relink}
AST rewrite $+$ \texttt{query/validate} execution and back-translation)
is $0.70 \pm 0.79$~s per query (median $0.85$~s, range $[0.00, 3.05]$~s),
on top of the baseline configuration's $1.16 \pm 0.37$~s LLM generation
(median $1.07$~s, range $[0.82, 3.38]$~s) that the MCP path reuses as
its input. The distribution is right-skewed: just under half of the queries
($49\,\%$) clear the middleware in under $0.05$~s; the tail comes from
the single LLM re-prompt issued when validation flags the relinked SQL.
The MCP-mediated end-to-end path is therefore $\approx 1.86$~s per query,
with a second LLM call incurred only on validation failure,
confirming the bounded-cost property posited in Section~III: at most
one initial generation plus one optional re-prompt per query.

Per-query cost was metered across every arm at August 2026 list
prices: single-call arms average $240$--$336$ input tokens
(USD $0.0004$--$0.0009$ per query); the agentic loop averages
$2507$--$4838$ tokens (USD $0.0061$--$0.0072$). A full two-database,
five-operator sweep costs under USD $3$ per model; the $199$-item
EvoSchema run cost USD $0.72$ (Haiku) and USD $1.25$ (GPT-4o). All middleware runs CPU-only, MiniLM validator included. With
Table~\ref{tab:scalability}'s sub-$5$\,ms primitives, the marginal
cost is the single optional re-prompt.

\subsection{Error Analysis of the Stale-Schema Baseline}

Table~\ref{tab:erroranalysis} splits the stale-schema baseline
outcomes by error category and operator. Each query is assigned one of
four verdicts: \emph{correct} (matches the gold result set),
\emph{execution\_error} (the SQL fails to run against the evolved schema),
\emph{silent\_failure} (the SQL runs but returns a wrong, non-empty result
that a naive check would accept), and \emph{wrong\_result} (the SQL runs
and returns a result that fails the gold-equality check). Execution errors
dominate every operator: $63$ of the $76$ baseline queries ($83\,\%$) fail
outright because the stale SQL references renamed or restructured objects
that no longer exist. The five \emph{wrong\_result} cases concentrate on
the merge and rename operators, where a stale identifier occasionally
resolves against a surviving object and yields a runnable but incorrect
query. No baseline query produced an undetected silent failure on this
database, which keeps the diff-and-relink pathway --- rather than
silent-failure detection --- as the dominant source of recovery in this
pilot.

\begin{table}[t]
\centering
\caption{Stale-schema baseline outcomes by error category, per operator
(Claude Haiku~4.5, \texttt{concert\_singer}). Corr.\ = correct;
Exec.\ = execution error; Silent = silent failure; Wrong = wrong result.}
\label{tab:erroranalysis}
\begin{tabular}{lccccc}
\toprule
\textbf{Operator} & \textbf{N} & \textbf{Corr.} & \textbf{Exec.} & \textbf{Silent} & \textbf{Wrong} \\
\midrule
TABLE\_RENAME  & 16 & 1 & 14 & 0 & 1 \\
TABLE\_SPLIT   & 16 & 1 & 14 & 0 & 1 \\
TABLE\_MERGE   & 14 & 2 & 12 & 0 & 0 \\
COLUMN\_RENAME & 16 & 1 & 14 & 0 & 1 \\
COLUMN\_MERGE  & 14 & 3 & 9  & 0 & 2 \\
\midrule
Total          & 76 & 8 & 63 & 0 & 5 \\
\bottomrule
\end{tabular}
\end{table}

\subsection{Ablation Study}

To attribute the end-to-end recovery to individual primitives, each MCP
component was disabled in turn while the other two were left intact, across
all five operators on both the \texttt{concert\_singer} and \texttt{hr\_1}
databases with Claude Haiku~4.5. Config~A is the full pipeline (control);
Config~B disables \texttt{schema/fingerprint} (the diff classifier, so
\texttt{query/relink} receives an empty event list); Config~C disables
\texttt{query/relink} (the stale SQL is passed straight to validate); and
Config~D disables \texttt{query/validate} (relinked SQL is returned with no
execution check and no LLM re-prompt). Table~\ref{tab:ablation} reports
pooled Execution Accuracy and the change against the full pipeline.

\begin{table}[t]
\centering
\caption{Ablation: pooled Execution Accuracy per database with one primitive
removed. $\Delta$ is the change against Config~A (Full MCP).}
\label{tab:ablation}
\begin{tabular}{llcc}
\toprule
\textbf{Database} & \textbf{Configuration} & \textbf{EX} & $\Delta$ \textbf{vs A} \\
\midrule
\texttt{concert\_singer} & A --- Full MCP       & 0.59 & ---     \\
\texttt{concert\_singer} & B --- No Fingerprint & 0.63 & $+0.04$ \\
\texttt{concert\_singer} & C --- No Relink      & 0.59 & $+0.00$ \\
\texttt{concert\_singer} & D --- No Validate    & 0.36 & $-0.24$ \\
\midrule
\texttt{hr\_1}           & A --- Full MCP       & 0.63 & ---     \\
\texttt{hr\_1}           & B --- No Fingerprint & 0.61 & $-0.03$ \\
\texttt{hr\_1}           & C --- No Relink      & 0.63 & $+0.00$ \\
\texttt{hr\_1}           & D --- No Validate    & 0.34 & $-0.29$ \\
\bottomrule
\end{tabular}
\end{table}

The ablation separates \emph{accuracy} benefits from \emph{systems}
benefits: latency, cost, reproducibility, interoperability.
Three findings stand out. First, removing \texttt{query/validate}
(Config~D) causes the largest drops --- between $-0.24$ and $-0.30$ pooled
--- concentrated on TABLE\_SPLIT, TABLE\_MERGE, and COLUMN\_MERGE. These
are precisely the operators the deterministic AST rewriter does not fully
handle, so they depend on the validate~$\rightarrow$~LLM-re-prompt fallback
to produce correct SQL; without validate, the pipeline ships the
unrecoverable stale SQL and every such query fails. Second, removing
\texttt{query/relink} (Config~C) yields a $0.000$ change everywhere under
Haiku~4.5: when the AST path is absent, validate still fires on the stale
SQL and the LLM re-prompt --- given the real diff --- recovers the same
queries the AST path would have handled. The relinker's contribution is
therefore a systems benefit, not an
accuracy benefit under a capable model: a deterministic
$\sim$$0.002$\,s AST pass that replaces an extra $\sim$$1.3$\,s LLM
call on the operators it covers, with bit-identical outputs across
repeated runs. The LLM re-prompt path guarantees neither.
The accuracy contribution concentrates in \texttt{query/validate}
(Config~D, $-0.24$ to $-0.29$), and detection in
\texttt{schema/fingerprint}, whose diff is what makes the recovery
targeted rather than blind regeneration.
Third, removing \texttt{schema/fingerprint} (Config~B) is mostly
neutral; the one anomaly is COLUMN\_MERGE on \texttt{concert\_singer}, where
Config~B ($0.50$) in fact outperforms Config~A ($0.21$), because a
placeholder diff string proves less misleading to the re-prompt than the
verbatim MERGE event the AST path injects. We frame this honestly: under a
capable model the deterministic relinker buys a fast, reproducible path
rather than headline accuracy, and whether its guarantees become more
valuable on weaker or cheaper models is left to the full evaluation.

\subsection{Cross-Model Generalisation of the MCP Server}

A central claim of this work is that the schema-aware recovery layer is
model-agnostic: the same Model Context Protocol server should restore
execution accuracy regardless of which pre-trained client generates the
initial SQL. The three primitives operate on artefacts, not on the
model. \texttt{schema/fingerprint} reads the database schema and emits
a classified diff; \texttt{query/relink} rewrites SQL through a
deterministic \texttt{sqlglot} abstract-syntax-tree pass driven by
that diff; \texttt{query/validate} executes the candidate query and
returns a three-state verdict. No step inspects the identity, weights,
or provider of the client that produced the SQL. The model is
therefore the independent variable and the server is held constant
across every configuration. The empirical runs reinforce this separation through
transport diversity: the four clients reach the harness over three distinct
integration paths --- Claude Haiku~4.5 through the native Anthropic SDK,
Gemini 2.5 Flash through the Google GenAI SDK, and both GPT-4o and
Llama~3.3~70B through an OpenAI-compatible HTTP endpoint --- yet the identical
fingerprint, relink, and validate pipeline is invoked in each case. Any
recovery observed across these clients is attributable to the server,
because the only component shared by all of them is the server itself.

Table~\ref{tab:crossmodel} reports Execution Accuracy for the stale,
refreshed, and MCP-mediated paths, together with the Recovery Rate, for
each client on both databases, with RR computed against the
refreshed-schema ceiling by Equation~\ref{eq:rr}.

\begin{table}[t]
\centering
\footnotesize
\caption{Cross-model execution accuracy and recovery rate (per database,
all five operators pooled).}
\label{tab:crossmodel}
\begin{tabular}{llcccc}
\toprule
\textbf{Model} & \textbf{Database} & \textbf{EX stale} & \textbf{EX refr.} & \textbf{EX MCP} & \textbf{RR} \\
\midrule
Haiku 4.5  & \texttt{concert\_singer} & 0.11 & 0.71 & 0.59 & 0.80 \\
Haiku 4.5  & \texttt{hr\_1}           & 0.21 & 0.67 & 0.62 & 0.89 \\
GPT-4o     & \texttt{concert\_singer} & 0.16 & 0.82 & 0.79 & 0.96 \\
GPT-4o     & \texttt{hr\_1}           & 0.21 & 0.68 & 0.66 & 0.94 \\
Gemini 2.5 Flash & \texttt{concert\_singer} & 0.16 & 0.88 & 0.83 & 0.93 \\
Gemini 2.5 Flash & \texttt{hr\_1}           & 0.21 & 0.68 & 0.66 & 0.94 \\
Llama 3.3 70B    & \texttt{concert\_singer} & 0.16 & 0.86 & 0.80 & 0.93 \\
Llama 3.3 70B    & \texttt{hr\_1}           & 0.22 & 0.72 & 0.71 & 0.97 \\
\bottomrule
\end{tabular}
\end{table}

All four clients completed the full evaluation matrix (both databases,
all five operators, 76~queries each). The MCP path raises Execution
Accuracy from near-zero on the stale schema to within a few points of
the refreshed-schema ceiling across every model. Recovery Rates by
client: Haiku~4.5 RR\,=\,0.80--0.89, GPT-4o RR\,=\,0.94--0.96,
Gemini~2.5 Flash RR\,=\,0.93--0.94, Llama~3.3~70B RR\,=\,0.93--0.97.
McNemar's paired test rejects the null hypothesis of no difference for
MCP versus the stale baseline on all eight model--database combinations
($b \geq 31$, $c = 0$; minimum $\chi^2 = 31$, $p < 0.001$), confirming
that the gain is not attributable to chance and is not explained by a
single error-correction retry.
How stable are these numbers under \texttt{temperature}=0 API
non-determinism? Each query was re-run five times (Haiku~4.5, both
databases, all operators). The stale, error-feedback, and MCP outcomes
were $100\,\%$ stable per query across all five seeds; agreement on
the pre-schema and refreshed arms was $0.993$ (two flips). A GPT-4o three-seed sweep on the same matrix shows the
stale, refreshed, and error-feedback outcomes $100\,\%$ stable with
MCP agreement $0.967$. The five flips, itemised in the repository,
isolate provider-side non-determinism on the re-prompt path. The
reported accuracies are stable point estimates, not single-draw
artefacts.

The AST rewriter can rebuild the affected identifiers or join
structure from the schema diff alone, so table-rename, table-split and
table-merge all recover strongly. The column-merge operator is the
consistent exception: when two columns are combined through an
expression, the original components are no longer recoverable from the
post-schema, the rewriter cannot invent the inverse transformation,
and Execution Accuracy on the MCP path stays at the stale level. This
boundary is a property of the operator rather than the client, and it
appears for every model evaluated. Three of the four models produced
RR above $1.0$ on at least one database (strongest: Gemini~2.5 Flash
on \texttt{concert\_singer} TABLE\_RENAME, RR\,=\,1.30). This occurs
when relinking a previously correct query yields higher accuracy than
regenerating it from scratch, since regeneration reintroduces the
model's own generation errors. The effect is real but small, and at
the present sample size the refreshed-schema ceiling carries
non-trivial sampling variance. Recovery Rate is therefore read as
recovery toward, and occasionally beyond, the refreshed-schema
reference rather than a strict upper bound; larger query sets in the
journal extension will tighten this estimate.

\subsection{MCP versus Vendor-Native Function Calling}
\label{sec:vendor}
Does the transport itself matter? The same three primitives were
re-exposed through vendor-native function calling
(Anthropic tool use for Claude Haiku~4.5, OpenAI function calling for
GPT-4o) in an agentic loop capped at eight tool turns. A third arm drops the
protocol entirely: the classified diff is pasted into a single
regeneration prompt (\emph{diff-in-prompt}). Every arm sees the same
middleware, byte for byte, and runs on both databases across all five
operators, $152$ scored queries per model. The re-run stale baseline matched the canonical baseline on every
query. Table~\ref{tab:vendor} shows the outcome. On
accuracy, the transports are statistically indistinguishable. Vendor-native
function calling reaches $\text{EX} = 0.658$ (Haiku) and $0.737$
(GPT-4o) against the MCP path's $0.605$ and $0.724$; McNemar's exact
test gives $p = 0.169$ and $p = 0.774$, Holm-corrected within each
model's test family. The diff-in-prompt arm matches vendor-native on
both models ($p = 1.00$ and $p = 0.774$). What carried the recovery was the classified schema diff itself; the
delivery protocol contributed nothing measurable. The design predicts
exactly this, since both surfaces route the same primitives over the
same payloads. Where the arms do differ is where the design says they
should: token spend and tool behaviour. The agentic vendor loop consumes
$9$--$14\times$ the tokens of the single-call arms: mean $4838$ vs.\
$336$ input tokens per query for Haiku, USD $0.0061$ vs.\ $0.0005$ per
query. The vendors also drove the primitives differently: Haiku invoked the
deterministic \texttt{relink\_sql} tool on $11\,\%$ of queries,
preferring to rewrite SQL itself from the diff; GPT-4o invoked it on
$86\,\%$. MCP's fixed server-side pipeline makes this asymmetry
irrelevant.

\begin{table}[t]
\centering
\footnotesize
\caption{MCP versus vendor-native function calling and diff-in-prompt
(pooled, both databases, five operators, $n=152$ per model). RR by
Equation~\ref{eq:rr} against the canonical refreshed ceiling. McNemar
$p$ is the exact two-sided test of the arm against the MCP path.}
\label{tab:vendor}
\begin{tabular}{llccc}
\toprule
\textbf{Model} & \textbf{Arm} & \textbf{EX} & \textbf{RR} & $p$ \textbf{vs MCP} \\
\midrule
Haiku 4.5 & MCP            & 0.605 & 0.840 & --- \\
Haiku 4.5 & Vendor-native  & 0.658 & 0.938 & 0.169 \\
Haiku 4.5 & Diff-in-prompt & 0.664 & 0.951 & 0.045 \\
GPT-4o    & MCP            & 0.724 & 0.953 & --- \\
GPT-4o    & Vendor-native  & 0.737 & 0.977 & 0.774 \\
GPT-4o    & Diff-in-prompt & 0.724 & 0.953 & 1.000 \\
\bottomrule
\end{tabular}
\end{table}

\subsection{Real-Benchmark Evaluation on EvoSchema}
\label{sec:evoschema}
The pilot query sets are authored; EvoSchema is not. To test the
mechanism on real benchmark data, the full four-configuration harness
was run on a deterministic, stratified subset of EvoSchema itself,
which perturbs the BIRD-dev databases. The subset samples $223$ items;
$199$ pass a gold-SQL validation gate requiring both gold queries to
execute with matching, non-empty results across the pre/post
databases. The $24$ exclusions are tabulated in the repository.
Coverage spans all eleven BIRD-dev databases and all five operator
classes. The five in-scope operator files total $5768$ items, so this
remains a subset evaluation. Table~\ref{tab:evoschema} reports the
results. The stale-schema baseline collapses almost completely on this
harder substrate: $\text{EX} = 0.000$ for Haiku and $0.020$ for
GPT-4o. Error feedback again fails to recover ($\leq 0.020$). The
MCP-mediated path restores $\text{EX}$ to $0.322$ (Haiku) and $0.412$
(GPT-4o), at or above the refreshed-schema ceiling of $0.296$ and
$0.362$ (RR $= 1.08$ and $1.15$). Both recoveries are significant at
$p < 10^{-6}$ (McNemar versus stale, Holm-corrected). Recovery Rates
above $1.0$ mean repair beats regeneration on harder questions, the
same effect as in Section~\ref{sec:results}-C, here on real benchmark
data. The absolute ceilings of $0.30$--$0.36$ reflect BIRD-dev
difficulty and bind every configuration alike; the mechanism closes
the floor-to-ceiling gap fully.

\begin{table}[t]
\centering
\footnotesize
\caption{EvoSchema subset (BIRD-dev substrate): $199$ gate-validated
items across all 11 databases and five operators. Err-fb =
error-feedback baseline.}
\label{tab:evoschema}
\begin{tabular}{lcccccc}
\toprule
\textbf{Model} & $n$ & \textbf{EX stale} & \textbf{EX refr.} & \textbf{EX err-fb} & \textbf{EX MCP} & \textbf{RR} \\
\midrule
Haiku 4.5 & 199 & 0.000 & 0.296 & 0.010 & 0.322 & 1.08 \\
GPT-4o    & 199 & 0.020 & 0.362 & 0.020 & 0.412 & 1.15 \\
\bottomrule
\end{tabular}
\end{table}

\subsection{Scalability}

To confirm the bounded-cost design at schema sizes beyond the pilot
databases, the \texttt{schema/fingerprint} diff and the \texttt{query/relink}
AST pass were timed on synthetic schemas of $10$, $100$, and $1000$ tables
($25$-iteration medians). Table~\ref{tab:scalability} reports the results.
Both primitives stay under $5$\,ms at $1000$ tables (fingerprint
median $3.4$\,ms, relink median $4.7$\,ms) and grow sub-linearly in
the table count.

\begin{table}[t]
\centering
\caption{Scalability of the fingerprint diff and AST relink passes by schema
size (median and IQR over 25 iterations, seconds).}
\label{tab:scalability}
\begin{tabular}{lcccc}
\toprule
\textbf{Tables} & \textbf{FP median} & \textbf{FP IQR} & \textbf{Relink median} & \textbf{Relink IQR} \\
\midrule
10   & 0.00005 & $<$0.00001 & 0.00020 & 0.00001 \\
100  & 0.00038 & 0.00001 & 0.00061 & 0.00002 \\
1000 & 0.00341 & 0.00014 & 0.00467 & 0.00031 \\
\bottomrule
\end{tabular}
\end{table}

% --------------------------------------------------------------- %
\section{Discussion}
% --------------------------------------------------------------- %

\paragraph*{Scope of the recovery-rate claim}
The pooled Recovery Rate of $0.80$ covers all five EvoSchema
operators on \texttt{concert\_singer} ($n = 76$). For the four straightforward rename
and merge operators, identifier substitution is correct by construction
--- every renamed identifier in a parsed query is mechanically
rewritten --- but the rewriter cannot normalise value-level predicate
mismatches the LLM introduces, as Q05 demonstrates. The column-merge
operator is the consistent exception (Table~\ref{tab:ablation}): when
two columns are combined through an expression, the original components
are no longer recoverable from the post-schema, so the rewriter cannot
invent the inverse transformation. These two failure modes ---
value-encoding mismatch and expression-based column merge ---
account for the $0.20$ residual. We therefore make no generalisation
claim to the silent-semantic failures that arise when a query returns
rows under a perturbed schema but answers the wrong question. The full-benchmark evaluation will extend the per-operator recovery
analysis beyond the subset of Section~\ref{sec:evoschema}.

\paragraph*{Isolation of schema retrieval as the independent variable}
The MCP-mediated configuration deliberately reuses the baseline
generation (the stale-schema LLM call) as its input rather than
re-prompting from scratch. This removes generation stochasticity as a
source of variance and ensures that any accuracy change is
attributable to the relink and validate primitives rather than to
lucky regeneration. The empirical consequence is that the achievable
MCP execution accuracy is bounded above by the refreshed-schema
ceiling: $\text{EX} = 0.59$ against $0.71$, a $0.12$ gap from
column-merge failures and the Q05 value-encoding mismatch, with no
other loss on the MCP path.

\paragraph*{Silent-failure detection}
The detector combines an empty-result-on-affirmative-question
check with a sentence-embedding cosine back-translation test
(Table~\ref{tab:sfd}). On the 60-case labelled set it attains
precision $1.00$ / recall $0.375$ / F1 $0.545$, versus $0.92$ / $0.30$
for the v1 token-overlap prototype. Precision comes first by design.
False positives are the only way \texttt{query/validate} can actively
harm accuracy, and the labelled evaluation shows it never suppresses a
faithful query. The remaining misses sit on mutations whose
back-translations stay semantically close to the original question
(wrong-column, dropped predicate); closing them needs
result-set--aware checks, not better similarity thresholds.

\paragraph*{MCP versus vendor-native tool use}
Section~\ref{sec:vendor} reports the like-for-like comparison. Both
surfaces route the same primitives, so accuracy comes out
statistically indistinguishable, and the diff-in-prompt arm shows the
classified diff carries the recovery on its own. If accuracy cannot
separate the two, what does? Three measured properties.
(a)~Integration cost: one server served every client unchanged, no
vendor adapter; tool-loop behaviour differed sharply (relink usage
$11\,\%$ vs.\ $86\,\%$) and the pipeline absorbed it. (b)~Client
cost: the agentic loop pays the overhead metered in
Section~\ref{sec:vendor}; under MCP the session Resource holds schema
state and the client pays one generation plus at most one re-prompt.
(c)~Audit surface: the Hou et al.~\cite{hou2025mcp} threat taxonomy
attaches to MCP's standardised protocol; vendor-native calling carries
the same risks without that audit point.

\paragraph*{Security implications for production deployment}
Exposing database access through any tool protocol enlarges the
attack surface, and MCP is no exception. Hou et al.~\cite{hou2025mcp}
identify a four-category attacker taxonomy across the MCP lifecycle;
tool poisoning is a prominent risk. Three properties of the proposed
server bound its exposure. The execution surface is read-only:
\texttt{query/validate} runs candidate SQL over a connection that
carries no DDL or DML privileges, so neither a compromised client nor
an injected prompt can mutate schema or data through the server. The
primitives are narrow and typed (a fingerprint resource, a rewrite of
caller-supplied SQL, a three-verdict validator) rather than a general
\texttt{run\_sql} tool. The server never composes new queries from natural language itself,
limiting prompt-to-SQL injection to the client's own capability. The Resource discloses schema metadata only, and fingerprints are
one-way BLAKE3 digests, so no data-exfiltration channel is added
beyond the query results already returned. Residual risks remain: registry-entry impersonation, result-set
leakage through the client context, and the \texttt{query/relink} LLM
fallback, which sends stale SQL and diffs to an external API and
should use a local model for sensitive schemas. A deployment checklist
ships with the reference implementation.

\paragraph*{Limitations and remediation paths}
Both residual failure modes have a concrete next action. Column
merges first. The post-schema alone cannot tell the rewriter how two
columns were combined, so the inverse expression stays out of reach.
The diff plus the original question can. We enriched the LLM-fallback
prompt with deterministic per-operator rewrite rules (an equality
filter over a merged column becomes a \texttt{LIKE} pattern over the
merge separator) and added one bounded retry for the case where an
affirmative question comes back empty. A hardening run built this way brought column-merge back to
near-ceiling on both pilot databases. Those figures post-date the
evaluated pilot configuration, so they live in the open-source
repository and will anchor the full evaluation. Q05-class
value-encoding mismatches are already detected by the empty-result
heuristic and the embedding validator (Table~\ref{tab:sfd}). What recovery still needs is the column's real value space in the
re-prompt, which the read-only \texttt{query/validate} connection can
supply at bounded cost.

\paragraph*{Threats to validity}
First, the pooled headline covers two Spider sub-domains
(\texttt{concert\_singer}, \texttt{hr\_1}) with authored query sets;
breadth across domains, query styles, and naturally occurring (rather
than synthesised) schema histories requires the full EvoSchema
benchmark beyond the subset of Section~\ref{sec:evoschema}.
Second, the evaluation uses SQLite; database engines
that enforce stricter type checking may reveal silent failures invisible
here. Third, the MCP-mediated configuration shares the baseline's
initial generation pass, which removes generation stochasticity as a
source of variance but also bounds the achievable ceiling at the
refreshed-schema reference accuracy.
Fourth, the silent-failure evaluation uses a 60-case labelled set
built from deterministic mutations of the pilot gold queries;
naturally occurring silent failures may distribute differently.
These constraints are acknowledged; the full evaluation addresses
them.

% --------------------------------------------------------------- %
\section{Conclusion and Future Work}
% --------------------------------------------------------------- %

This paper presented the first design and pilot evaluation, to our
knowledge, of a Model-Context-Protocol-mediated approach to schema
evolution in NL-to-SQL. The proposed server exposes three runtime
primitives whose composition enables an off-the-shelf LLM client to detect,
represent, and recover from schema change without fine-tuning. On a
76-query pooled evaluation over a perturbed Spider
\texttt{concert\_singer} database spanning all five EvoSchema operator
classes, the MCP-mediated configuration achieves Recovery Rate $= 0.80$
($\text{EX}_{\text{mcp}} = 0.59$ vs.\ stale-schema baseline $0.11$ and
refreshed-schema ceiling $0.71$), with the $0.20$ residual
attributable to column-merge expression reconstruction and
value-encoding mismatches the identifier-substitution rewriter cannot
normalise.
Three further results support the design: vendor-native function
calling matches MCP accuracy; the mechanism transfers to
EvoSchema/BIRD-dev (RR $> 1$); the validator reaches precision $1.00$
/ F1 $0.545$.

Future work will extend the evaluation to the full Spider, BIRD, and
EvoSchema benchmarks, with per-operator recovery rates and McNemar's
analysis across all splits. Three pipeline items follow from the failure analysis: raise
silent-failure recall without giving up zero false positives, supply
value-space samples in the re-prompt so Q05-class detections become
recoveries, and promote the column-merge fallback into the default
pipeline. The reference implementation
is released under an MIT licence so that subsequent research on
MCP-mediated database access can build on a shared standardised
substrate rather than ad-hoc vendor adapters.

% --------------------------------------------------------------- %
\section*{Acknowledgments}
% --------------------------------------------------------------- %

The authors would like to thank the anonymous reviewers for their suggestions to improve the paper.

% --------------------------------------------------------------- %
% References [1]-[28], canonical numbering mirrored from build_docx.py.
\begin{thebibliography}{99}

\bibitem{vaswani2017transformer}
A.~Vaswani \textit{et~al.}, ``Attention is all you need,'' in
\emph{Proc.\ NeurIPS}, 2017.

\bibitem{zhong2017seq2sql}
V.~Zhong, C.~Xiong, and R.~Socher, ``Seq2SQL: Generating structured queries
from natural language using reinforcement learning,'' arXiv:1709.00103,
2017.

\bibitem{li2023resdsql}
H.~Li, J.~Zhang, C.~Li, and H.~Chen, ``RESDSQL: Decoupling schema linking
and skeleton parsing for text-to-SQL,'' in \emph{Proc.\ AAAI}, 2023.

\bibitem{pourreza2023dinsql}
M.~Pourreza and D.~Rafiei, ``DIN-SQL: Decomposed in-context learning of
text-to-SQL with self-correction,'' in \emph{Proc.\ NeurIPS}, 2023.

\bibitem{wang2025macsql}
B.~Wang \textit{et~al.}, ``MAC-SQL: A multi-agent collaborative framework
for text-to-SQL,'' in \emph{Proc.\ COLING}, 2025.

\bibitem{pourreza2025chasesql}
M.~Pourreza \textit{et~al.}, ``CHASE-SQL: Multi-path reasoning and
preference optimised candidate selection for text-to-SQL,'' in
\emph{Proc.\ ICLR}, 2025.

\bibitem{talaei2024chess}
S.~Talaei, M.~Pourreza, Y.-C.~Chang, A.~Mirhoseini, and A.~Saberi,
``CHESS: Contextual harnessing for efficient SQL synthesis,''
arXiv:2405.16755, 2024. [Preprint; not peer-reviewed.]

\bibitem{yu2018spider}
T.~Yu \textit{et~al.}, ``Spider: A large-scale human-labeled dataset for
complex and cross-domain semantic parsing and text-to-SQL task,'' in
\emph{Proc.\ EMNLP}, 2018.

\bibitem{li2023bird}
J.~Li \textit{et~al.}, ``Can LLM already serve as a database interface? A
big bench for large-scale database grounded text-to-SQLs,'' in
\emph{Proc.\ NeurIPS}, 2023.

\bibitem{brahmia2024schemaevol}
Z.~Brahmia, F.~Grandi, and B.~Oliboni, ``A literature review on schema
evolution in databases,'' \emph{Computing Open}, vol.~02, 2024,
doi:10.1142/s2972370124300012.

\bibitem{zhang2025evoschema}
T.~Zhang \textit{et~al.}, ``EvoSchema: Towards text-to-SQL robustness
against schema evolution,'' \emph{Proc.\ VLDB Endowment}, vol.~18, no.~10,
pp.~3655--3668, 2025, doi:10.14778/3748191.3748222.

\bibitem{anthropic2024mcp}
Anthropic, ``Model Context Protocol: Specification (revision 2024-11-05),''
2024. [Online]. Available: \url{https://modelcontextprotocol.io}

\bibitem{hou2025mcp}
X.~Hou, Y.~Zhao, S.~Wang, and H.~Wang, ``Model Context Protocol (MCP):
Landscape, security threats, and future research directions,''
arXiv:2503.23278, 2025. [Preprint; not peer-reviewed.]

\bibitem{agater2026toolbelt}
J.~Agater, L.~Herradi, M.~Mimouni, A.~Memari, and J.~M.~G\'omez,
``Toolbelt-MCP: Exploring the Model Context Protocol for tool-use of
large language models utilizing graphs with relational databases,'' in
\emph{Artificial Intelligence XLII} (Lecture Notes in Computer Science,
vol.~16302), M.~Bramer and F.~Stahl, Eds.\ Cham, Switzerland: Springer,
2026, pp.~472--477, doi:10.1007/978-3-032-11442-6\_38.

\bibitem{masayoshi2025ehrmcp}
K.~Masayoshi \textit{et~al.}, ``EHR-MCP: Real-world evaluation of clinical
information retrieval by large language models via Model Context
Protocol,'' arXiv:2509.15957, 2025. [Preprint; not peer-reviewed.]

\bibitem{shi2025survey}
L.~Shi, Z.~Tang, N.~Zhang, X.~Zhang, and Z.~Yang, ``A survey on employing
large language models for text-to-SQL tasks,'' \emph{ACM Comput.\ Surv.},
vol.~58, no.~2, Art.~54, 2025, doi:10.1145/3737873.

\bibitem{hong2025tkde}
Z.~Hong, Z.~Yuan, Q.~Zhang, H.~Chen, J.~Dong, F.~Huang, and X.~Huang,
``Next-Generation Database Interfaces: A Survey of LLM-Based Text-to-SQL,''
\emph{IEEE Trans.\ Knowl.\ Data Eng.}, vol.~37, no.~12, pp.~7328--7345, 2025,
doi:10.1109/TKDE.2025.3609486.

\bibitem{li2026deepeyesql}
B.~Li, C.~Chen, Z.~Xue, Y.~Mei, and Y.~Luo, ``DeepEye-SQL: A
software-engineering-inspired text-to-SQL framework,''
\emph{Proc.\ ACM Manag.\ Data}, vol.~4, no.~3, Art.~158, 2026,
doi:10.1145/3788284.

\bibitem{lei2026crl}
J.~Lei, R.~He, and Y.~Wang, ``Enhanced Text-to-SQL using cooperative
reinforcement learning for small language models,'' \emph{Neurocomputing},
vol.~673, art.~132865, 2026, doi:10.1016/j.neucom.2026.132865.

\bibitem{ojuri2025agents}
S.~Ojuri, T.~A.~Han, R.~Chiong, and A.~Di~Stefano, ``Optimizing
text-to-SQL conversion techniques through the integration of intelligent
agents and large language models,'' \emph{Inf.\ Process.\ Manag.},
vol.~62, p.~104136, 2025, doi:10.1016/j.ipm.2025.104136.

\bibitem{nascimento2026kgsql}
E.~R.~Nascimento \textit{et~al.}, ``A text-to-SQL strategy based on large
language models and knowledge graphs for real-world databases,''
\emph{Data Knowl.\ Eng.}, vol.~164, p.~102580, 2026,
doi:10.1016/j.datak.2026.102580.

\bibitem{liao2026infosql}
H.~Liao, S.~Chen, Y.~Xiao, L.~Xiang, and F.~Min, ``Large-scale
text-to-SQL generation with adversarial defense,'' \emph{Information
Sciences}, vol.~733, art.~122942, 2026, doi:10.1016/j.ins.2025.122942.

\bibitem{chen2025gbvsql}
D.~Chen, X.~Wang, S.~Ren, Q.~Ma, P.~Zhao, and A.~Liu, ``GBV-SQL: Guided
generation and SQL2Text back-translation validation for multi-agent
Text2SQL,'' arXiv:2509.12612, 2025. [Preprint; not peer-reviewed.]

\bibitem{li2025cmsql}
X.~Li, J.~You, H.~Li, J.~Peng, X.~Chen, and Z.~Guo, ``CM-SQL: A
cross-model consistency framework for text-to-SQL,'' \emph{Neurocomputing},
vol.~658, art.~131708, 2025, doi:10.1016/j.neucom.2025.131708.

\bibitem{jian2026trustsql}
A.~Jian \textit{et~al.}, ``TRUST-SQL: Tool-integrated multi-turn
reinforcement learning for text-to-SQL over unknown schemas,''
arXiv preprint arXiv:2603.16448, 2026. [Preprint; not peer-reviewed.]

\bibitem{lei2026srspsql}
J.~Lei, Z.~Li, and Y.~Wang, ``SRSPSQL: A dual-stage Text-to-SQL framework
with semantic rewriting and schema pruning,'' \emph{Information and Software
Technology}, vol.~194, art.~108064, 2026, doi:10.1016/j.infsof.2026.108064.

\bibitem{avro2024fp}
Apache Software Foundation, ``Apache Avro 1.12.0 Specification: Schema
Fingerprints,'' 2024. [Online]. Available:
\url{https://avro.apache.org/docs/1.12.0/specification/#schema-fingerprints}

\bibitem{phan2025askdb}
X.-Q.~Phan, T.-H.~Mai, T.-D.~Dinh, M.-T.~Nguyen, and L.-S.~L\^e, ``AskDB:
An LLM agent for natural language interaction with relational databases,''
arXiv:2511.16131, 2025. [Preprint; not peer-reviewed.]

\bibitem{bogdanov2026agentnlq}
O.~Bogdanov, Y.~Jung, C.~Dhir, P.~Gaddam, S.~Jain, L.~Tumati,
V.~Parthasarathy, and A.~Shirgaonkar, ``AgentNLQ: A general-purpose agent
for natural language to SQL,'' arXiv:2605.19010, 2026. [Preprint; not
peer-reviewed.]

\bibitem{hassini2025dynaquery}
A.~Hassini, ``DynaQuery: A self-adapting framework for querying structured
and multimodal data,'' arXiv:2510.18029, 2025. [Preprint; not
peer-reviewed.]

\end{thebibliography}

\end{document}

